% !PRINTED 775
\weberpage{775}
Hieraus folgt die Unzulässigkeit unserer Annahme, dass $f_n(x)$ nach der Formel (11) im Körper $R(\alpha)$ in zwei Factoren zerlegbar ist.

\begin{center}
\rule{4cm}{0.4pt}
\end{center}

Damit lässt sich nun aufs Neue allgemein beweisen, dass die Gleichung, deren Wurzeln die sämmtlichen primitiven $m^{\text{ten}}$ Einheitswurzeln sind, deren Grad $\varphi(m)$ ist, auch wenn $m$ eine beliebige zusammengesetzte Zahl ist, irreducibel ist. Dies ist mit folgendem Satze gleichbedeutend:

\gesperrt{Wenn eine ganze Function $\varphi(x)$ mit rationalen Coëfficienten verschwindet, wenn für $x$ eine primitive $m^{\text{te}}$ Einheitswurzel gesetzt wird, so verschwindet $\varphi(x)$ auch, wenn für $x$ irgend eine andere primitive $m^{\text{te}}$ Einheitswurzel gesetzt wird.}

Wenn $m$ eine Primzahlpotenz ist, so ist der Satz in dem Vorhergehenden schon bewiesen. Wir wenden daher die vollständige Induction an, nehmen ihn für $m^{\text{te}}$ Einheitswurzeln als erwiesen an und leiten ihn für $mn^{\text{te}}$ Einheitswurzeln her, wenn $n = p^k$ und $p$ eine in $m$ nicht aufgehende Primzahl ist.

Jede $mn^{\text{te}}$ Einheitswurzel $\beta$ kann in der Form

\[
\beta = \alpha r
\]

dargestellt werden, worin $\alpha$ eine $m^{\text{te}}$, $r$ eine $n^{\text{te}}$ Einheitswurzel ist. Was zu beweisen ist, lässt sich so aussprechen, dass, wenn $\beta$ die Wurzel einer rationalen Gleichung $\Phi(x) = 1$ ist, % [sic]
auch alle anderen primitiven Einheitswurzeln vom Grade $mn$ Wurzeln dieser Gleichung sein müssen.

Ist, wie oben, $f_n(x) = 0$ die Gleichung, deren Wurzeln die Grössen $r$ sind, so ist, wie wir bewiesen haben, $f_n(x)$ nicht in rationale Factoren zerlegbar, auch nicht in solche, die rational von $\alpha$ abhängen. Die Function $f_n(\alpha^{-1}x)$, die für $x = \beta$ verschwindet, kann also auch nicht in rationale Factoren zerlegt werden, weil aus einer Zerlegung der Form

\[
f_n(\alpha^{-1}x) = \varphi(x,\alpha)\ \psi(x,\alpha)
\]

folgen würde:

\[
f_n(x) = \varphi(\alpha x,\alpha)\ \psi(\alpha x,\alpha),
\]

was nicht bestehen kann. Da hiernach $f_n(\alpha^{-1}x)$ eine Wurzel
%CONT>>
